\documentclass{article}


% ============================================================================
% NeurIPS 2026 style. Using [preprint] so real author names are shown
% (not anonymized) and the footer reads "Preprint. Work in progress."
% instead of the double-blind submission notice. Switch to [final] once
% the course confirms a camera-ready-style render is wanted, or drop the
% option entirely to see the anonymized double-blind version with line
% numbers.
% ============================================================================
\usepackage[preprint]{neurips_2026}

\usepackage[utf8]{inputenc}
\usepackage[T1]{fontenc}
\usepackage{hyperref}
\usepackage{url}
\usepackage{booktabs}
\usepackage{amsfonts}
\usepackage{amsmath}
\usepackage{amssymb}
\usepackage{nicefrac}
\usepackage{microtype}
\usepackage{xcolor}
\usepackage{graphicx}

% ---------------------------------------------------------------------------
% Placeholder helpers: everything wrapped in \pending{...} is a stand-in for
% content that depends on the 900-step training runs, which are still in
% progress on the lab/rented GPU as of this draft. Search for "PENDING" to
% find every spot that still needs real numbers before submission.
% ---------------------------------------------------------------------------
\newcommand{\pending}[1]{\textcolor{red!70!black}{\bf [PENDING: #1]}}
\newcommand{\figplaceholder}[2]{%
  \begin{center}
    \fbox{\parbox{0.92\linewidth}{\centering\vspace{2.2cm}
      \textcolor{red!70!black}{\bf Figure placeholder}\\[2pt]
      \small #2
      \vspace{2.2cm}}}
  \end{center}
  \captionof{figure}{#1}
}
\usepackage{caption}

\title{Entropy-Regularized Rewards for Self-Rewarded Training: \\
Delaying Reward Hacking in Majority-Vote RL for Mathematical Reasoning}

% TODO: confirm affiliation / department wording and emails before submitting.
\author{
  Daniel Shavit \\
  \texttt{TODO: Department, University} \\
  \texttt{shavit.daniel.2@gmail.com} \\
  \And
  Yuval Zohar \\
  \texttt{TODO: Department, University} \\
  \texttt{TODO: email} \\
}

\begin{document}

\maketitle

\begin{abstract}
Self-Rewarded Training (SRT) \citep{shafayat2025srt} trains large reasoning
models on mathematical problem-solving without ground-truth labels, using
majority-vote agreement across sampled rollouts as a proxy reward. SRT
improves accuracy early in training, but the original paper shows that
extended training causes the proxy reward to be gamed: the model learns to
maximize self-agreement without becoming more correct, and performance
collapses. We propose adding a per-rollout entropy (surprisal) bonus to the
SRT reward -- a classical exploration-preserving regularizer from
policy-gradient RL \citep{williams1991function,mnih2016a3c,haarnoja2018sac} --
to test whether it delays or reduces this collapse. We implement a
lightweight, dependency-minimal replication of SRT's RLOO training loop on
Qwen2.5-Math-1.5B over the DAPO dataset \citep{yu2025dapo}, and compare four
reward conditions -- plain SRT ($\alpha{=}0$), two entropy-bonus strengths
($\alpha{=}0.3$, $\alpha{=}0.5$), and a ground-truth oracle ceiling -- each
trained for 900 steps with periodic held-out evaluation. \pending{final
900-step accuracy/gap numbers and headline result sentence -- runs are still
executing as of this draft}. A preliminary frozen-model pilot across three
off-the-shelf models confirms that the agreement-accuracy gap SRT relies on
is a real, measurable phenomenon prior to any training, motivating the
intervention studied here.
\end{abstract}

% =============================================================================
\section{Introduction}
% =============================================================================

Reinforcement learning has become the dominant recipe for turning a
pretrained language model into a strong mathematical reasoner, but almost
every successful recipe -- from RLHF to verifiable-reward RL for math --
depends on a reward signal that requires either human-labeled preferences or
a ground-truth answer key. Both are expensive to obtain at scale, which
motivates \emph{label-free} RL: can a model improve at reasoning using only
signals it can compute about its own outputs, with no external verifier?

\citet{shafayat2025srt} study exactly this question with \textbf{Self-Rewarded
Training (SRT)}. At every RL step, SRT samples several candidate solutions to
the same problem from the \emph{current} policy, takes their majority-vote
answer as a stand-in for ground truth, and rewards each rollout for agreeing
with that majority. This builds directly on self-consistency decoding
\citep{wang2022selfconsistency}, which showed that majority voting over
multiple chain-of-thought samples is itself a strong test-time proxy for
correctness. SRT turns that proxy into a training signal: no labels are
needed because the model is, in effect, both student and teacher.

This works, at first. \citeauthor{shafayat2025srt} report genuine accuracy
gains with no ground-truth supervision at all. But they also identify the
method's central failure mode, which motivates this project: with
\textbf{prolonged training}, the policy can learn to maximize the
self-consistency reward without becoming more \emph{correct} -- it collapses
onto confidently repeating one answer per question, right or wrong, because
that guarantees perfect agreement with itself. The self-consistency reward is
maximized precisely as true accuracy stagnates or falls. This is a textbook
instance of reward hacking (or Goodhart's law: a proxy stops being a good
measure once it is directly optimized), and the original paper leaves
\emph{fixing} it as an open problem rather than a solved one.

\paragraph{Our contribution.} We propose a direct, mechanistically-motivated
fix: augment SRT's reward with a small bonus proportional to how
\emph{surprising} (rare) a rollout's answer is relative to the other sampled
rollouts for that question. This is an application of entropy regularization
-- a decades-old technique for preventing premature convergence in
policy-gradient RL \citep{williams1991function}, popularized by A3C
\citep{mnih2016a3c} and formalized as an explicit training objective by
maximum-entropy RL methods such as Soft Actor-Critic
\citep{haarnoja2018sac} -- applied specifically to SRT's self-consistency
collapse, which the original paper does not attempt to address. We implement
this as a lightweight, from-scratch RLOO \citep{kool2019buy,ahmadian2024back}
training loop, add a ground-truth oracle condition as an upper-bound
reference the original paper does not report, and run a controlled
comparison across four reward conditions at 900 training steps each -- three
times the budget of our initial 300-step feasibility run, chosen because
reward hacking in the original paper's own experiments only becomes visible
after substantial training.

The rest of this report is organized as follows. Section~\ref{sec:background}
reviews the RL background needed to follow SRT and our modification.
Section~\ref{sec:contribution} states our reward function and hypotheses
precisely. Section~\ref{sec:methodology} describes the experimental setup.
Section~\ref{sec:results} reports results (with the main training comparison
pending completion of the 900-step runs). Section~\ref{sec:discussion}
discusses limitations and future work.


% =============================================================================
\section{Background}
\label{sec:background}
% =============================================================================

\subsection{Policy-gradient RL and RLOO}

We treat generating a solution as a sequence of token-level actions under a
policy $\pi_\theta$ (the language model). Policy-gradient methods update
$\theta$ in the direction of $\mathbb{E}[A \cdot \nabla_\theta \log
\pi_\theta(\text{sequence})]$, where $A$ is an advantage estimate: how much
better a sampled sequence's reward was than some baseline. \textbf{RLOO}
(REINFORCE Leave-One-Out) \citep{kool2019buy,ahmadian2024back} is a
critic-free advantage estimator well suited to this setting: given $k$
rollouts $\{r_1,\dots,r_k\}$ for the same prompt with rewards
$\{R_1,\dots,R_k\}$, the baseline for rollout $i$ is the mean reward of the
\emph{other} $k-1$ rollouts,
\[
A_i = R_i - \frac{1}{k-1}\sum_{j\neq i} R_j ,
\]
avoiding the need to train a separate value network (unlike PPO) while
remaining unbiased. SRT and our replication both use RLOO.

\subsection{Self-consistency and Self-Rewarded Training}

\citet{wang2022selfconsistency} observed that sampling multiple
chain-of-thought solutions to the same problem and taking a majority vote
over their final answers improves accuracy at inference time, without any
additional training -- correct reasoning paths tend to agree with each
other more than incorrect ones. \textbf{SRT} \citep{shafayat2025srt} turns
this inference-time heuristic into a \emph{training} signal: at every RL
step, it samples $k$ rollouts for a training prompt, computes their majority
answer, and assigns reward $R_i = \mathbf{1}[a_i = \text{majority}(a_1,\dots,a_k)]$
to each rollout $i$ -- 1 if it agrees with the crowd, 0 otherwise. Because the
``teacher'' (the majority vote) is recomputed from the model's own,
constantly-updating outputs at every step, SRT differs from other
label-free methods that instead fix a majority-vote snapshot from a frozen
model and train against it once (e.g.\ SFT-on-majority-vote, DPO, or
ScPO-style preference optimization). This continuously-updating property is
precisely what lets SRT track a moving target more closely than its
fixed-snapshot cousins -- and, as discussed next, is also why it is
uniquely vulnerable to drifting into confidently-wrong self-agreement.

\subsection{Reward hacking via self-consistency collapse}

Because $R_i$ only ever measures agreement with the model's own majority
vote, and never checks the real answer, nothing in the objective
distinguishes ``confidently correct'' from ``confidently wrong.'' If the
policy's answer distribution for a given question narrows -- even for
reasons unrelated to correctness -- the majority answer becomes easier to
match, reward goes up, and the update reinforces that narrowing further. This
is a feedback loop: overconfidence is self-reinforcing under a
purely-agreement-based reward. \citeauthor{shafayat2025srt} document exactly
this collapse empirically: after enough training, the self-consistency
reward keeps climbing (near its ceiling) while true accuracy on held-out
questions stagnates or drops. We track this with the same diagnostic used in
the original paper's own analysis: the \textbf{agreement-accuracy gap},
\[
\text{gap} = \text{agreement} - \text{accuracy},
\]
where agreement is the fraction of a question's $k$ rollouts that match the
majority answer (needs no ground truth), and accuracy is whether that
majority answer is actually correct (needs ground truth, computed only for
diagnostic/evaluation purposes, never fed into the SRT training reward). A
gap near zero means the model's confidence tracks its correctness; a large,
growing gap is the reward-hacking signature.

\subsection{Entropy regularization in policy-gradient RL}

Encouraging a policy to keep some residual randomness, rather than
collapsing deterministically onto one action, is a long-standing fix for
premature convergence in policy-gradient methods. \citet{williams1991function}
is among the earliest to use an entropy-like regularization term in
connectionist RL. \citet{mnih2016a3c}'s A3C algorithm adds the policy's
entropy directly to the training objective, scaled by a coefficient, to
discourage collapsing onto one action before exploration is complete --
structurally the same shape as the bonus we add here.
\citet{haarnoja2018sac}'s Soft Actor-Critic generalizes this into
\emph{maximum-entropy RL}: optimizing for reward \emph{plus} entropy as a
first-class part of the objective, rather than a minor add-on. Separately,
\citet{pathak2017curiosity}'s curiosity-driven exploration rewards an agent
for encountering low-probability (\emph{surprising}) outcomes, which is the
more direct ancestor of the per-rollout \emph{surprisal} bonus we use (as
opposed to a batch-level entropy term). None of this machinery is new to RL
in general; what SRT's own paper does not do is apply it to its own
self-consistency collapse, which is the gap this project targets.


% =============================================================================
\section{Contribution}
\label{sec:contribution}
% =============================================================================

\subsection{Reward function}

For a training question with $k$ rollout answers $a_1,\dots,a_k$ (some
possibly unparseable / \texttt{None}), let $p(a)$ be the empirical fraction
of \emph{valid} rollouts among the $k$ that produced answer $a$. We define
the per-rollout reward
\begin{equation}
  R_i \;=\; \underbrace{\mathbf{1}[a_i = \text{majority}(a_1,\dots,a_k)]}_{\text{SRT's self-consistency reward}}
  \;+\; \alpha \cdot \underbrace{\big(-\log p(a_i)\big)}_{\text{surprisal bonus}},
  \label{eq:reward}
\end{equation}
with $\alpha \geq 0$ a fixed coefficient controlling the bonus strength.
$\alpha = 0$ recovers plain SRT exactly. Rollouts with no parseable answer
receive $R_i = 0$ and are excluded from the $p(a)$ estimate. Concretely: if 5
of 8 rollouts answer ``84'', 2 answer ``91'', and 1 answers ``72'', then
$p(84){=}0.625$, $p(91){=}0.25$, $p(72){=}0.125$, giving surprisal values of
approximately $0.47$, $1.39$, and $2.08$ respectively -- so at $\alpha{=}0.5$
the lone, rare ``72'' rollout earns almost as much reward ($1.04$) as the
five majority ``84'' rollouts ($1.0$ each), instead of the $0$ it would get
under plain SRT. This is the mechanism by which the bonus is meant to keep
disagreement worthwhile enough that the policy does not fully collapse onto
one answer before it has had the chance to discover the correct one.

\subsection{Why this specifically targets SRT's failure mode}

Under plain SRT ($\alpha{=}0$), any rollout that disagrees with the current
majority gets exactly zero reward, regardless of whether it happens to be
correct. As the policy's output distribution narrows over training, this
structurally starves rare (but possibly correct) answers of any gradient
signal, accelerating collapse. The surprisal term in
Eq.~\ref{eq:reward} restores a reward for disagreement, scaled by how rare
the disagreement currently is -- so the more the policy narrows, the larger
the bonus for stepping outside the majority becomes, directly counteracting
the runaway-confidence feedback loop described in
Section~\ref{sec:background}. Because SRT's majority-vote target keeps
evolving with the model (Section~\ref{sec:background}), it is more exposed to
this loop than fixed-snapshot alternatives (SFT/DPO/ScPO on a single
majority-vote pass); an exploration-preserving regularizer is a natural
candidate fix for that specific dynamic. We are explicit that the entropy
bonus \emph{mechanism} is not new (Section~\ref{sec:background}); the
contribution is diagnosing that SRT's particular training dynamic is exactly
the kind of premature-convergence failure this class of regularizer was
designed for, and testing that empirically.

\subsection{Oracle ceiling}

We additionally implement a \textbf{ground-truth oracle} condition,
$R_i = \mathbf{1}[a_i = \text{ground truth}]$, with no majority vote and no
entropy bonus at all -- the classic verifiable-reward RL setup used when a
real answer key \emph{is} available. The original SRT paper does not report
this comparison. We include it as an upper-bound reference: it uses the same
RLOO training loop and the same rollout budget as every other condition, so
its accuracy curve shows what is achievable if the reward signal were
perfect, letting us judge how much of the accuracy gap between SRT variants
and ``ideal'' RL is attributable to the majority-vote proxy itself versus to
reward hacking specifically.

\subsection{Hypotheses}

\begin{itemize}
  \item \textbf{H1 (collapse under budget):} given enough training steps,
  plain SRT ($\alpha{=}0$) will show a growing agreement-accuracy gap,
  replicating the original paper's collapse finding at our (much smaller)
  model scale.
  \item \textbf{H2 (bonus delays collapse):} $\alpha > 0$ conditions will
  show a smaller and/or later-growing gap than $\alpha{=}0$, at the possible
  cost of slower or lower peak accuracy early in training (since reward is
  partly diverted away from pure majority-matching).
\end{itemize}
Both hypotheses are evaluated in Section~\ref{sec:results} once the 900-step
runs complete; our earlier 300-step feasibility run (Section~\ref{sec:methodology})
did not run long enough to distinguish between the conditions on either axis,
which is precisely why the training budget was extended for the results
reported here.


% =============================================================================
\section{Methodology}
\label{sec:methodology}
% =============================================================================

\subsection{Task, data, and model}

We use the \textbf{DAPO} dataset curation \citep{yu2025dapo}, specifically
the unlabeled variant \texttt{ftajwar/deduplicated\_dapo\_dataset} used by
the original SRT codebase. During training, the visible ``ground truth''
field is a placeholder (\texttt{LABEL\_BY\_SELF\_CONSISTENCY}) by design --
no real answer is available to the training loop except in the oracle
condition, where the real answer is read from a separate,
deliberately-segregated field
(\texttt{reward\_model.solution\_hidden\_during\_training}) that is never
exposed to the non-oracle reward computation. Evaluation uses a disjoint,
held-out 100-question test set (fixed \texttt{random\_state=42}); a nested,
fixed 20-question subset of it is used for periodic during-training checks
so that trend curves compare the same questions at every checkpoint, and the
full 100 questions are used for a single, final evaluation pass at the end
of each run. Training rows overlapping with the eval set are excluded to
avoid leakage.

We use \textbf{Qwen2.5-Math-1.5B} as the base policy, loaded via
\texttt{transformers} in fp16. This is a much smaller model than typically
used in SRT-style RL papers, chosen to fit the project's compute budget
(a memory-constrained $\sim$11GB lab GPU, later supplemented with a rented
GPU for the longer runs reported here).

\subsection{Training loop}

We implement SRT's algorithm directly (RLOO advantage, majority-vote /
entropy-bonus / oracle rewards) in a single, dependency-minimal script rather
than using the original paper's full \texttt{verl}+vLLM+FSDP training stack,
which pulls in heavyweight distributed-training dependencies not needed at
our scale. Answer extraction (\verb|\boxed{...}| parsing) and
ground-truth equivalence checking reuse the original SRT repository's own
\texttt{math\_verify}-based scorer directly (via a thin dynamic-import
wrapper), rather than a reimplementation, so that e.g.\ \verb|\boxed{0.5}|
and \verb|\boxed{1/2}| are scored as equal exactly as they would be in the
original codebase, keeping our numbers comparable to it.

Each training step: sample one training prompt; generate $k{=}8$ rollouts at
temperature $1.0$ (rollouts that are truncated before an end-of-sequence
token are automatically regenerated once at a longer token budget); extract
and score answers; compute rewards per Eq.~\ref{eq:reward} (or the oracle
reward); compute RLOO advantages; and take one SGD step (learning rate
$10^{-6}$, momentum $0.9$, gradient norm clipped to $1.0$) using gradients
accumulated one rollout at a time to fit the vocabulary-projection forward
pass in memory. Before any real training run, we run an automated
\textbf{gradient-invariance sanity check} that verifies the entropy bonus
actually changes the computed RLOO advantages on a synthetic example
(catching a ``zero-gradient bug'' class of implementation error before
burning GPU time on a run that would silently reduce to plain SRT regardless
of $\alpha$). Long runs checkpoint after every step and support
crash-resume, which mattered in practice on a shared/rented GPU.

\subsection{Experimental conditions}

We compare four conditions, each trained for \textbf{900 steps}: $\alpha{=}0.0$
(plain SRT, the paper's method), $\alpha{=}0.3$, $\alpha{=}0.5$ (our primary
proposed configuration), and the ground-truth oracle. Held-out evaluation
(20-question subset) runs every 15 steps to produce accuracy/gap/entropy
trend curves; a single final evaluation over all 100 test questions is run
at the end of each condition's 900 steps for the headline comparison number.
900 steps was chosen as three times our initial 300-step feasibility run's
budget: at 300 steps we observed no measurable gap growth in any condition
(\pending{300-step numbers omitted from this report by design -- see
Section~\ref{sec:results} for why}), consistent with the original paper's
own observation that reward-hacking collapse is a phenomenon of
\emph{extended} training, not early training. 900 steps is our best estimate,
within the project's compute budget, of a horizon long enough to observe the
onset of collapse if it occurs at this model scale.

\subsection{Preliminary pilot: does the gap exist before any training?}

Before committing compute to full RL training runs, we first ran a
\textbf{frozen-model pilot}: no training at all, just the same
majority-vote-vs.-ground-truth measurement applied to several off-the-shelf
models' zero-shot behavior on a fixed 20-question subset, to confirm the
agreement-accuracy gap this project studies is a real, measurable phenomenon
worth building a training intervention around. Models tested: Qwen2.5-32B
(lab GPU), Llama-4-Scout-17B (Groq API), and OpenAI's o4-mini. Results are
reported in Section~\ref{sec:results}; this pilot is complete and is not
affected by the still-running 900-step training experiments.

\subsection{Metrics}

For every evaluated question, agreement is the fraction of $k$ rollouts
matching the majority answer; accuracy (majority-vote accuracy) is whether
that majority answer matches the real answer; gap is agreement minus
accuracy (Section~\ref{sec:background}); mean entropy is the Shannon entropy
of the answer distribution across the $k$ rollouts. Loss is logged but,
following the original paper's convention, is not used to compare
conditions -- it reflects the policy-gradient objective's internal scale,
not task performance.


% =============================================================================
\section{Experiments and Results}
\label{sec:results}
% =============================================================================

\subsection{Preliminary pilot: frozen-model agreement-accuracy gap}

Table~\ref{tab:pilot} reports the completed frozen-model pilot
(Section~\ref{sec:methodology}): each model's zero-shot accuracy, 8-rollout
majority-vote accuracy, and agreement, all on the same fixed evaluation
subset, no training involved.

\begin{table}[h]
  \caption{Frozen-model pilot: majority-vote accuracy and agreement, no
  training. Confirms the agreement-accuracy gap is present, and varies in
  size, across off-the-shelf models before any RL intervention.}
  \label{tab:pilot}
  \centering
  \begin{tabular}{lccc}
    \toprule
    Model & Majority-vote accuracy & Agreement & Gap \\
    \midrule
    o4-mini                        & 0.95 & 0.91 & $-0.04$ \\
    Llama-4-Scout-17B-16E-Instruct & 0.40 & 0.68 & $+0.28$ \\
    Qwen2.5-32B                    & 0.45 & 0.53 & $+0.08$ \\
    \bottomrule
  \end{tabular}
\end{table}

The gap varies substantially by model: o4-mini's agreement already tracks
its (high) accuracy closely, while Llama-4-Scout shows a much larger gap --
its rollouts agree with each other noticeably more often than they are
actually correct, i.e.\ exactly the pattern SRT's reward would reinforce if
this model were trained with it. This motivates studying the collapse
dynamic directly through training, rather than assuming it is negligible at
the model scale used in our main experiment.

\subsection{Main training comparison (900 steps)}

\pending{This subsection reports the core result of the report and is not
yet complete. As of this draft, the $\alpha{=}0.0$ run is in progress on the
lab server and the $\alpha{=}0.3$ run needs to be relaunched after an
earlier interruption; $\alpha{=}0.5$ and the oracle condition are queued
behind them. The table, figure, and paragraph below are structured to be
filled in directly once all four 900-step runs finish their final
100-question evaluation -- only the numbers/curves need to be dropped in,
the surrounding analysis text should be revisited but the shape should not
need to change.}

\begin{table}[h]
  \caption{Final (100-question) evaluation after 900 training steps, by
  condition. \pending{fill in once runs complete}}
  \label{tab:main-results}
  \centering
  \begin{tabular}{lccc}
    \toprule
    Condition & Test accuracy & Agreement & Gap \\
    \midrule
    $\alpha = 0.0$ (plain SRT)     & \pending{TBD} & \pending{TBD} & \pending{TBD} \\
    $\alpha = 0.3$                  & \pending{TBD} & \pending{TBD} & \pending{TBD} \\
    $\alpha = 0.5$ (main method)    & \pending{TBD} & \pending{TBD} & \pending{TBD} \\
    Oracle (ceiling)                & \pending{TBD} & \pending{TBD} & \pending{TBD} \\
    \bottomrule
  \end{tabular}
\end{table}

\figplaceholder{Test accuracy (left axis convention) and agreement-accuracy
gap (right axis convention) vs.\ training step, for all four conditions,
evaluated every 15 steps on the fixed 20-question subset.}{To be generated
from \texttt{lab\_alpha0.0.json}, \texttt{lab\_alpha0.3.json},
\texttt{lab\_alpha0.5.json}, and \texttt{lab\_oracle.json} once the 900-step
runs finish -- one accuracy curve and one gap curve per condition, sharing a
step axis.}

\pending{Result narrative: state here whether H1 held (did $\alpha{=}0.0$'s
gap grow over the 900 steps, unlike at 300 steps?) and whether H2 held (did
$\alpha{=}0.3$/$0.5$ show a smaller or later-growing gap than $\alpha{=}0.0$,
and at what cost, if any, to accuracy?). Compare all three non-oracle
conditions' final accuracy against the oracle ceiling to quantify how much
of the achievable accuracy is left on the table by the label-free reward
proxies.}


% =============================================================================
\section{Discussion}
\label{sec:discussion}
% =============================================================================

\subsection{Interpreting the result}
\pending{Write this once Section~\ref{sec:results}'s table is filled in.
Two shapes this could take: (a) if collapse appears in $\alpha{=}0.0$ by
step 900 and is measurably smaller in $\alpha{=}0.3$/$0.5$, this section
should discuss the accuracy/robustness trade-off and which $\alpha$ appears
best-justified; (b) if no condition shows collapse even at 900 steps, this
section should discuss what that implies about the training-budget scale at
which SRT's failure mode emerges at this model size, and whether smaller
models are simply slower to reach the confidence levels needed for the
majority-vote reward to start being gamed.}

\subsection{Limitations}

Several limitations apply regardless of the final 900-step numbers.
\textbf{Model scale}: Qwen2.5-Math-1.5B is far smaller than the models used
in the original SRT paper and in most RL-for-math literature; conclusions
here may not transfer to larger models, which could reach the confidence
levels needed for collapse faster or slower than observed here.
\textbf{Single seed per condition}: compute constraints meant each of the
four 900-step conditions was run once, with no repeated seeds -- reported
numbers have no error bars or significance testing, and single noisy
checkpoints could be mistaken for genuine trend changes.
\textbf{Rollout budget}: $k{=}8$ rollouts per question is smaller than some
published SRT-style setups use, which weakens both the majority-vote signal
itself and the resolution of the surprisal estimate $p(a)$.
\textbf{Domain}: all training and evaluation is on DAPO math problems; we
make no claim that the entropy bonus's effect (if any) generalizes to other
verifiable-reward domains (code, logic puzzles) where SRT-style training has
also been applied.
\textbf{Fixed $\alpha$}: we test two discrete $\alpha$ values rather than a
continuous sweep or a schedule (e.g.\ annealing $\alpha$ down over training),
so we cannot rule out that some untested $\alpha$, or a time-varying
schedule, performs meaningfully better than either value tested here.

\subsection{Future work}

Natural extensions include: repeating the 900-step comparison with multiple
seeds to get error bars on the gap curves; sweeping $\alpha$ more finely (or
annealing it) rather than comparing two fixed values; testing at a larger
model scale closer to the original paper's; and testing whether an entropy
bonus tuned on math transfers to other SRT application domains. If collapse
is confirmed at 900 steps for $\alpha{=}0$, a natural next question is
whether even longer training eventually causes collapse under the bonus too,
just delayed -- i.e.\ whether the bonus changes the location of the failure
mode or removes it.

% =============================================================================
% References
% =============================================================================
\begin{thebibliography}{99}

\bibitem[Ahmadian et al.(2024)]{ahmadian2024back}
Arash Ahmadian, Chris Cremer, Matthias Gall\'{e}, Marzieh Fadaee, Julia
Kreutzer, Olivier Pietquin, Ahmet \"{U}st\"{u}n, and Sara Hooker.
\newblock Back to basics: Revisiting REINFORCE-style optimization for
learning from human feedback in LLMs.
\newblock \emph{arXiv preprint arXiv:2402.14740}, 2024.

\bibitem[Haarnoja et al.(2018)]{haarnoja2018sac}
Tuomas Haarnoja, Aurick Zhou, Pieter Abbeel, and Sergey Levine.
\newblock Soft actor-critic: Off-policy maximum entropy deep reinforcement
learning with a stochastic actor.
\newblock In \emph{Proceedings of the 35th International Conference on
Machine Learning (ICML)}, 2018.

\bibitem[Kool et al.(2019)]{kool2019buy}
Wouter Kool, Herke van Hoof, and Max Welling.
\newblock Buy 4 REINFORCE samples, get a baseline for free!
\newblock In \emph{ICLR 2019 Workshop DeepRLStructPred}, 2019.

\bibitem[Mnih et al.(2016)]{mnih2016a3c}
Volodymyr Mnih, Adri\`{a} Puigdom\`{e}nech Badia, Mehdi Mirza, Alex Graves,
Tim Harley, Timothy~P. Lillicrap, David Silver, and Koray Kavukcuoglu.
\newblock Asynchronous methods for deep reinforcement learning.
\newblock In \emph{Proceedings of the 33rd International Conference on
Machine Learning (ICML)}, 2016.

\bibitem[Pathak et al.(2017)]{pathak2017curiosity}
Deepak Pathak, Pulkit Agrawal, Alexei~A. Efros, and Trevor Darrell.
\newblock Curiosity-driven exploration by self-supervised prediction.
\newblock In \emph{Proceedings of the 34th International Conference on
Machine Learning (ICML)}, 2017.

\bibitem[Shafayat et al.(2025)]{shafayat2025srt}
Sheikh Shafayat et al.
\newblock Can large reasoning models self-train?
\newblock \emph{arXiv preprint arXiv:2505.21444}, 2025.

\bibitem[Wang et al.(2022)]{wang2022selfconsistency}
Xuezhi Wang, Jason Wei, Dale Schuurmans, Quoc Le, Ed~Chi, Sharan Narang,
Aakanksha Chowdhery, and Denny Zhou.
\newblock Self-consistency improves chain of thought reasoning in language
models.
\newblock \emph{arXiv preprint arXiv:2203.11171}, 2022.

\bibitem[Williams and Peng(1991)]{williams1991function}
Ronald~J. Williams and Jing Peng.
\newblock Function optimization using connectionist reinforcement learning
algorithms.
\newblock \emph{Connection Science}, 3(3):241--268, 1991.

\bibitem[Yu et al.(2025)]{yu2025dapo}
Qiying Yu, Zheng Zhang, Ruofei Zhu, et al.
\newblock DAPO: An open-source LLM reinforcement learning system at scale.
\newblock \emph{arXiv preprint arXiv:2503.14476}, 2025.

\end{thebibliography}


% =============================================================================
\newpage
\appendix
\section*{NeurIPS Paper Checklist}
% =============================================================================

\begin{enumerate}

\item \textbf{Claims}\\
    Question: Do the main claims made in the abstract and introduction
    accurately reflect the paper's contributions and scope?\\
    Answer: \answerTODO{}\\
    Justification: \pending{Revisit once Section~\ref{sec:results} is
    filled in -- the abstract currently states the result is pending, which
    must be updated to match whatever H1/H2 outcome is found.}

\item \textbf{Limitations}\\
    Question: Does the paper discuss the limitations of the work performed
    by the authors?\\
    Answer: \answerYes{}\\
    Justification: Section~\ref{sec:discussion} discusses model scale, single-seed
    runs, rollout budget, domain restriction, and the fixed (non-swept)
    $\alpha$ values tested.

\item \textbf{Theory assumptions and proofs}\\
    Question: For each theoretical result, does the paper provide the full
    set of assumptions and a complete (and correct) proof?\\
    Answer: \answerNA{}\\
    Justification: The paper is empirical; it introduces no formal
    theorems.

\item \textbf{Experimental result reproducibility}\\
    Question: Does the paper fully disclose all the information needed to
    reproduce the main experimental results?\\
    Answer: \answerYes{}\\
    Justification: Section~\ref{sec:methodology} specifies the model,
    dataset, reward formula (Eq.~\ref{eq:reward}), rollout count, learning
    rate, optimizer, and evaluation protocol (fixed seeds, held-out
    question subsets) needed to reproduce the setup.

\item \textbf{Open access to data and code}\\
    Question: Does the paper provide open access to the data and code, with
    sufficient instructions to faithfully reproduce the main experimental
    results?\\
    Answer: \answerYes{}\\
    Justification: Code and logged results are available at
    \url{https://github.com/yuval-4241/RL_PROJECT}; the DAPO dataset used is
    the public \texttt{ftajwar/deduplicated\_dapo\_dataset}.

\item \textbf{Experimental setting/details}\\
    Question: Does the paper specify all the training and test details
    necessary to understand the results?\\
    Answer: \answerYes{}\\
    Justification: See Section~\ref{sec:methodology} (rollout count,
    temperature, learning rate, optimizer, gradient clipping, evaluation
    cadence, train/test split and leakage exclusion).

\item \textbf{Experiment statistical significance}\\
    Question: Does the paper report error bars or other appropriate
    statistical-significance information?\\
    Answer: \answerNo{}\\
    Justification: Each of the four 900-step conditions was run with a
    single seed due to compute/time constraints (Section~\ref{sec:discussion}); no
    repeated-seed error bars are reported. This is flagged as a limitation.

\item \textbf{Experiments compute resources}\\
    Question: Does the paper provide sufficient information on the compute
    resources needed to reproduce the experiments?\\
    Answer: \answerYes{}\\
    Justification: Training used a memory-constrained ($\sim$11GB) lab GPU
    plus a rented GPU for the longer 900-step runs
    (Section~\ref{sec:methodology}); \pending{insert exact GPU model and
    total wall-clock/GPU-hours once the 900-step runs finish}.

\item \textbf{Code of ethics}\\
    Question: Does the research conform with the NeurIPS Code of Ethics?\\
    Answer: \answerYes{}\\
    Justification: The project trains on a public, deduplicated math
    dataset and evaluates a reward-function modification; it poses no
    identified conflicts with the Code of Ethics.

\item \textbf{Broader impacts}\\
    Question: Does the paper discuss both potential positive and negative
    societal impacts?\\
    Answer: \answerNo{}\\
    Justification: This is a small-scale, class-project-level replication
    and extension of an existing published method on a public math dataset;
    we do not identify a direct path to negative applications distinct from
    the original SRT paper's own discussion.

\item \textbf{Safeguards}\\
    Question: Does the paper describe safeguards for responsible release of
    high-misuse-risk data or models?\\
    Answer: \answerNA{}\\
    Justification: No new model checkpoints or datasets are released
    publicly beyond code and logged metrics; the base model (Qwen2.5-Math-1.5B)
    is an existing open-weight release.

\item \textbf{Licenses for existing assets}\\
    Question: Are creators/owners of existing assets properly credited,
    with licenses and terms of use respected?\\
    Answer: \answerYes{}\\
    Justification: The base model (Qwen2.5-Math-1.5B), the DAPO dataset
    curation \citep{yu2025dapo}, and the original SRT codebase
    \citep{shafayat2025srt} are all cited; all are used under their
    respective public release terms.

\item \textbf{New assets}\\
    Question: Are new assets introduced in the paper well documented?\\
    Answer: \answerNA{}\\
    Justification: No new datasets or model checkpoints are introduced as
    standalone released assets; the training code is documented in the
    linked repository.

\item \textbf{Crowdsourcing and research with human subjects}\\
    Question: For crowdsourcing/human-subjects experiments, does the paper
    include full instructions and compensation details?\\
    Answer: \answerNA{}\\
    Justification: No human subjects or crowdsourcing were involved.

\item \textbf{IRB approvals}\\
    Question: Does the paper describe IRB approval (or equivalent) for
    human-subjects research?\\
    Answer: \answerNA{}\\
    Justification: No human-subjects research was conducted.

\item \textbf{Declaration of LLM usage}\\
    Question: Does the paper describe LLM usage if it is an important,
    original, or non-standard component of the core methods?\\
    Answer: \answerNo{}\\
    Justification: LLMs (Qwen2.5-Math-1.5B, and separately Qwen2.5-32B /
    Llama-4-Scout / o4-mini in the frozen-model pilot) are the
    \emph{subject} of the research, not a tool used to produce the core
    method or results; \pending{confirm and disclose here if any LLM
    assistance was used for code scaffolding, writing, or analysis during
    the project, per course/NeurIPS LLM-usage policy}.

\end{enumerate}

\end{document}
